\scanpage{24}
\margin{Consulter Weil, \unsure{Critères} d'\unsure{Équivalence|Équivalences}}

\uline{Théorie de l'équivalence d'Albanese des cycles}

1. \uline{L'anneau \deleted{\unsure{de l'|des}}des \inserted{classes d'}\unsure{équivalence|équivalences} d'Albanese}

\uline{Définition} 1.1. Soit $X$ un schéma \deleted{[projectif]} lisse sur le corps $k$, $z$ un cycle sur $X$. On dit que $z$ est alb-équiv. à $0$,

\margin{?}

si on peut trouver un schéma [projectif] lisse \deleted{\unsure{irréductible}}\inserted{connexe} $\deleted{\unsure{Y}}T$ sur $k$, \deleted{\unsure{et}}\deleted{\unsure{et}} un $0$-cycle \deleted{de degré $0$ $\unsure{\mathfrak a|\alpha|\mathfrak{t}}$} $\mathfrak a$ sur $\deleted{\unsure{Y}}T$, et un cycle $Z$ sur $X\times T$, tels que

(i) \inserted{$\mathfrak a$ est de degré $0$ et} $S_Y(\mathfrak a)=0$\qquad (ii) $Z(\mathfrak a)$ est défini et égal à $z$.

\uline{Proposition} 1.2. (i) Les cycles alb-équiv. à $0$ sur $X$ forment un sous-groupe \inserted{gradué $N^*(X)$} de $\mathrm{Cyc}^*(X)$, contenant le groupe des cycles lin. équiv. à $0$, de sorte que la connaissance de $N^*(X)$ équivaut à celle de son image $N'{}^*(X)$ dans \deleted{Cyc} $A^*(X)$ (l'anneau de Chow de $X$, défini par les classes à équiv. linéaire près), et $\mathrm{Cyc}(X)/N(X)\simeq\deleted{\unsure{\mathrm{Cyc}}}A^*(X)/N'{}^*(X)$.

% archive sheet number: 23
